\subsubsection{GloVe Training Objective}\label{sec:glove-training-objective}

We have already seen that GloVe starts from the global co-occurrence matrix. For every pair of words $i$ and $j$, we know $X_{ij}$, the number of times word $j$ appears in the context of word $i$. From this we compute:
\begin{equation*}
    P\left(j \mid i\right) = \frac{X_{ij}}{X_i}
\end{equation*}
Which expresses the probability of observing word $j$ near word $i$. The key question is now: \emph{\textbf{how do we transform these probabilities into word vectors?}}

\highspace
Suppose we have the words ``\emph{ice}'', ``\emph{steam}'', ``\emph{solid}'', and ``\emph{gas}''. From the corpus we observe:
\begin{center}
    \begin{tabular}{@{} c c @{}}
        \toprule
        \textbf{Pair} & \textbf{Co-occurrence} \\
        \midrule
        \emph{ice} - \emph{solid} & high \\
        \emph{steam} - \emph{gas} & high \\
        \emph{ice} - \emph{gas} & low \\
        \emph{steam} - \emph{solid} & low \\
        \bottomrule
    \end{tabular}
\end{center}
\noindent
Ideally we would like the vectors to satisfy ``\emph{ice}'' close to ``\emph{solid}'', ``\emph{steam}'' close to ``\emph{gas}'', without explicitly programming these relationships. Instead, the vectors themselves should explain the observed statistics.

\highspace
\textcolor{Green3}{\faIcon{question-circle} \textbf{What Should the Vectors Predict?}} The previous section introduced an important idea. The \textbf{meaning} of a word is reflected by how often it appears together with other words. Therefore GloVe asks: \emph{\textbf{can we choose vectors whose relationships reproduce the observed co-occurrence information?}} Unlike Word2Vec, which predicts neighboring words through a neural network, GloVe directly learns vectors that fit the global co-occurrence statistics.

\highspace
\begin{flushleft}
    \textcolor{Green3}{\faIcon{link} \textbf{Linear Relationship Between Vectors}}
\end{flushleft}
Imagine two words \emph{ice} and \emph{steam}, and one context word \emph{solid}. Suppose:
\begin{equation*}
    P\left(\emph{solid} \mid \emph{ice}\right) = 0.30
\end{equation*}
While:
\begin{equation*}
    P\left(\emph{solid} \mid \emph{steam}\right) = 0.01
\end{equation*}
Clearly, \emph{solid} is much more informative for \emph{ice} than for \emph{steam}. So the learned vectors should somehow encode this large difference. The question becomes: \emph{\textbf{what mathematical relationship between vectors best captures these probability differences?}}

\highspace
The GloVe authors observed that a very simple model works remarkably well. They assume that the \textbf{dot product} between two vectors should approximate the logarithm of their co-occurrence count. Specifically:
\begin{equation}
    w_i^T \tilde{w}_j + b_i + \tilde{b}_j \approx \log \left(X_{ij}\right)
\end{equation}
Where:
\begin{itemize}
    \item $w_i$ is the embedding of the target word (transposed for the dot product),
    
    
    \item $\tilde{w}_j$ is the embedding of the context word,
    
    
    \item $b_i, \tilde{b}_j$ are bias terms for the target and context words, respectively.
    
    
    \item $X_{ij}$ is the observed co-occurrence count of word $j$ in the context of word $i$.
    
    
    \item[\textcolor{Green3}{\faIcon{question-circle}}] \textcolor{Green3}{\textbf{Why the Dot Product?}} Remember what a dot product measures. If \hl{two vectors point in similar directions}:
    \begin{equation*}
        w_i^T \tilde{w}_{j}
    \end{equation*}
    Is \hl{large}. If they are unrelated, it is small. Therefore the dot product naturally measures how strongly two words are related.
    
    
    \item[\textcolor{Green3}{\faIcon{question-circle}}] \textcolor{Green3}{\textbf{Why the Logarithm?}} This is probably the most interesting part. Suppose we have:
    \begin{center}
        \begin{tabular}{@{} c c @{}}
            \toprule
            \textbf{Pair} & \textbf{Count} \\
            \midrule
            \emph{the} - \emph{of} & 1.000.000 \\
            \emph{cat} - \emph{animal} & 100 \\
            \emph{cat} - \emph{keyboard} & 1 \\
            \bottomrule
        \end{tabular}
    \end{center}
    Using raw counts is problematic because they vary enormously. A very common word like ``\emph{the}'' would dominate the optimization. Instead, GloVe compresses the scale:
    \begin{equation*}
        \log(1) = 0 \quad \log(100) \approx 4.6 \quad \log(1'000'000) \approx 13.8
    \end{equation*}
    The \hl{logarithm preserves the ordering while greatly reducing the range of values.} As a result, frequent word pairs no longer overwhelm the learning process.
\end{itemize}
This equation is the heart of GloVe.

\highspace
\begin{flushleft}
    \textcolor{Green3}{\faIcon{bullseye} \textbf{Learning the Embeddings}}
\end{flushleft}
Now we know what we want. For every observed pair:
\begin{equation*}
    w_i^T \tilde{w}_j + b_i + \tilde{b}_j
\end{equation*}
Should be as close as possible (approximately) to:
\begin{equation*}
    \log \left(X_{ij}\right)
\end{equation*}
Of course this \hl{equality cannot hold perfectly for every word pair simultaneously.} Instead, we \textbf{minimize the overall error across the entire corpus}.

\highspace
The GloVe authors define the following cost function:
\begin{equation}
    J = \sum_{i, j} f\left(X_{ij}\right) \cdot \left(w_i^T \tilde{w}_j + b_i + \tilde{b}_j - \log \left(X_{ij}\right)\right)^{2}
\end{equation}
Called the \textbf{weighted least squares objective}, it sums the squared differences between the predicted and actual log co-occurrence counts. The terms are:
\begin{itemize}
    \item \important{The Prediction}
    \begin{equation*}
        w_i^T \tilde{w}_j + b_i + \tilde{b}_j
    \end{equation*}
    This is the model's estimate of the relationship between two words.
    \item \important{The Target}
    \begin{equation*}
        \log \left(X_{ij}\right)
    \end{equation*}
    This comes directly from the corpus. It is the quantity we want to reproduce.
    \item \important{The Squared Error}
    \begin{equation*}
        \left(\text{prediction} - \text{target}\right)^{2}
    \end{equation*}
    If the prediction is accurate, the error is close to zero. Otherwise it becomes large.
    \item \important{The Weight Function}
    \begin{equation*}
        f\left(X_{ij}\right)
    \end{equation*}
    This is a very clever addition. Not all co-occurrence counts are equally reliable. For example:
    \begin{center}
        \begin{tabular}{@{} c c @{}}
            \toprule
            \textbf{Count} & \textbf{Reliability} \\
            \midrule
            1 & very noisy \\
            3 & noisy \\
            500 & reliable \\
            10.000 & reliable enough \\
            \bottomrule
        \end{tabular}
    \end{center}
    The weighting function assigns \textbf{less importance to rare co-occur-\break rences}, which are often due to chance, while giving more importance to frequent, reliable observations. At the same time, it \textbf{prevents extremely frequent pairs from dominating the optimization}.
\end{itemize}
In summary, GloVe learns word vectors by minimizing the difference between the \textbf{similarity predicted by the vectors} (their dot product plus biases) and the \textbf{similarity observed in the corpus} (the logarithm of the co-occurrence count), giving greater importance to reliable co-occurrence statistics than to noisy ones.

\highspace
In other words, \hl{GloVe learns one vector for every word such that the dot products between word vectors reproduce the co-occurrence statistics observed in the corpus.}